\chapter{Exchange Economies}\label{ch:b2:exchange-economies}

In 1962 a group of students traded an imaginary good in a classroom. Each buyer held a card with the
% ledger: b2:smith1962-t1
highest price she could pay, each seller a card with the lowest price he could accept, and anyone could
call out a bid or an offer at any moment. No one knew the others' cards, and no one computed an
equilibrium. Yet the predicted price, the one at which the supply built from the sellers' cards met the
demand built from the buyers', was \usd{2.00} in the first market, and the average contract price went from
\usd{1.80} in the first trading period to \usd{2.03} in the fifth. The chart shows six of the markets Smith ran:
each line is the average contract price of a trading period, as a percentage of the predicted price.

\begin{openingchart}{Average contract price in each trading period, as a percentage of the competitive
  price predicted from the buyers' and sellers' reservation prices, in six laboratory markets (tests 1,
  4A, 6A, 7, 9A and 10). The dashed line is the prediction.}{Smith (1962), Journal of Political
  Economy, Table~1}
\begin{tikzpicture}
\begin{axis}[omrecord, grid=none, ymajorgrids, xmin=0.7, xmax=6.3, ymin=40, ymax=120, xtick={1,...,6},
    xlabel={Trading period}, ylabel={Average price, per cent of predicted}, unbounded coords=jump,
    width=0.8\linewidth, height=0.5\linewidth, legend columns=3,
    legend style={at={(0.5,-0.22)}, anchor=north, font=\footnotesize}]
  \addplot[omDef, thick, mark=*] table[x=period, y=T1_ratio]
    {figdata/bachelor-2/06-exchange-economies/out/smith_wide.csv};
  \addplot[omThm, thick, mark=square*] table[x=period, y=T4A_ratio]
    {figdata/bachelor-2/06-exchange-economies/out/smith_wide.csv};
  \addplot[omExo, thick, mark=triangle*] table[x=period, y=T6A_ratio]
    {figdata/bachelor-2/06-exchange-economies/out/smith_wide.csv};
  \addplot[black, thick, mark=diamond*] table[x=period, y=T7_ratio]
    {figdata/bachelor-2/06-exchange-economies/out/smith_wide.csv};
  \addplot[black!50, thick, mark=o] table[x=period, y=T9A_ratio]
    {figdata/bachelor-2/06-exchange-economies/out/smith_wide.csv};
  \addplot[omDef!50, thick, densely dotted, mark=x] table[x=period, y=T10_ratio]
    {figdata/bachelor-2/06-exchange-economies/out/smith_wide.csv};
  \legend{Test 1, Test 4A, Test 6A, Test 7, Test 9A, Test 10}
  \draw[black!70, dashed] (axis cs:0.7,100) -- (axis cs:6.3,100);
\end{axis}
\end{tikzpicture}
\end{openingchart}

% figdata: smith_summary (t6a_first 49.2, t6a_last 101.4, t4a_last 107.1, t9a_last 90.3, share_alpha_fell 100)
Some markets opened far from the prediction: in test 6A the first period's average price was 49.2 per
cent of it, and by the fourth 101.4 per cent. Others stayed off: test 4A ended 7.1 per cent above, test
9A 9.7 per cent below. But in all ten tests that ran for three periods or more, prices clustered closer to
the prediction in the last period than in the first. Those were single markets, studied with the
partial-equilibrium tools of the Year 1 volume. An economy is many markets at once: what a household
supplies in one pays for what it demands in another, and every price affects every demand. This chapter
builds the simplest such economy, in which people only trade what they already have. It asks which
\omterm{def:b2:exchange-economies:allocation}{allocations} cannot be improved on, which prices clear every market at once, and whether such prices
exist, are unique and can be found by trading.

\begin{recall}
From the Year 1 volume: partial equilibrium in one market; consumer choice with a budget line, the
marginal rate of substitution and the tangency condition; Cobb--Douglas demands, which spend a fixed
share of income on each good. From \cref{ch:b2:insurance-and-risk-sharing}: Pareto efficiency of risk
sharing.
\end{recall}

\begin{notation}[Symbols of this chapter]\label{not:b2:exchange-economies:symbols}
Consumers are $i = A, B$ (or $i = 1, \dots, I$); goods $\ell = 1, 2$ (or $1, \dots, n$). Consumer $i$ has
the \omterm{def:b2:exchange-economies:endowment}{endowment} $\omega_i = (\omega_{i1}, \omega_{i2})$ and consumes $x_i = (x_{i1}, x_{i2})$; $\bar\omega
= \omega_A + \omega_B$ is the total \omterm{def:b2:exchange-economies:endowment}{endowment}. $u_i$ is a utility function and $\mathit{MRS}_i$ its
marginal rate of substitution of good 2 for good 1, $(\partial u_i/\partial x_{i1})/(\partial
u_i/\partial x_{i2})$. Prices are $p_1, p_2$; with good 2 as \omterm{def:b2:exchange-economies:numeraire}{numeraire}, $p = p_1/p_2$. $z_\ell(p)$ is the
excess demand for good $\ell$. With Cobb--Douglas utilities, $a$ and $b$ are the shares of income that
$A$ and $B$ spend on good 1.
\end{notation}

\section{The economy}\label{sec:b2:exchange-economies:economy}

\begin{definition}[Exchange economy]\label{def:b2:exchange-economies:exchange-economy}
An \emph{exchange economy}\index{exchange economy} consists of $I$ consumers and $n$ goods; each
consumer has preferences over bundles $x_i \in \mathbb{R}^n_{+}$ and an \omterm{def:b2:exchange-economies:endowment}{endowment}. There is no
production: the goods available are those the consumers hold.
\end{definition}

\begin{definition}[Endowment]\label{def:b2:exchange-economies:endowment}
A consumer's \emph{endowment}\index{endowment} $\omega_i \in \mathbb{R}^n_{+}$ is the bundle she holds
before any trade.
\end{definition}

\begin{definition}[Allocation]\label{def:b2:exchange-economies:allocation}
An \emph{allocation}\index{allocation} is a list $x = (x_1, \dots, x_I)$ of bundles, one per consumer.
\end{definition}

\begin{definition}[Feasible allocation]\label{def:b2:exchange-economies:feasible}
An \omterm{def:b2:exchange-economies:allocation}{allocation} is \emph{feasible}\index{feasible allocation} if $\sum_i x_i \le \sum_i \omega_i$, good by
good: the economy does not consume more of any good than it has. It is \emph{non-wasteful} if equality
holds.
\end{definition}

With two consumers and two goods, all non-wasteful \omterm{def:b2:exchange-economies:allocation}{allocations} fit in one picture.

\begin{definition}[Edgeworth box]\label{def:b2:exchange-economies:edgeworth-box}
The \emph{Edgeworth box}\index{Edgeworth box} of a two-consumer, two-good economy is the rectangle of
width $\bar\omega_1$ and height $\bar\omega_2$. A point in it is a non-wasteful \omterm{def:b2:exchange-economies:allocation}{allocation}: measured from
the bottom-left corner it is $A$'s bundle, measured from the top-right corner $B$'s.
\end{definition}

$A$'s indifference curves are drawn as usual from $A$'s origin; $B$'s are drawn from $B$'s origin, so
they curve the other way and $B$ prefers points nearer the bottom-left. A budget line through the
\omterm{def:b2:exchange-economies:endowment}{endowment} with slope $-p_1/p_2$ is a budget line for both consumers at once: every point on it is worth
as much as the \omterm{def:b2:exchange-economies:endowment}{endowment}, for $A$ and therefore for $B$.

The running example of this chapter has $\bar\omega = (10, 10)$, $\omega_A = (8, 2)$ and $\omega_B =
(2, 8)$: $A$ is rich in good 1 and $B$ in good 2. Utilities are $u_A = \tfrac12\ln x_{A1} + \tfrac12\ln
x_{A2}$ and $u_B = \tfrac23\ln x_{B1} + \tfrac13\ln x_{B2}$: $A$ spends half her income on each good, $B$
two thirds of his on good 1.

\begin{omfigure}
\begin{tikzpicture}[econ]
\begin{axis}[omrecord, grid=none, axis lines=box, xmin=0, xmax=10, ymin=0, ymax=10, xtick={0,2,4,6,8,10},
    ytick={0,2,4,6,8,10}, xlabel={Good 1 held by $A$}, ylabel={Good 2 held by $A$},
    width=0.6\linewidth, height=0.6\linewidth, clip=true]
  \addplot[omDef, thick, domain=3.05:10, samples=60] {30.4511/x};
  \addplot[omThm, thick, domain=0:9.2, samples=60] {10 - 98.9315/(10 - x)^2};
  \addplot[black, thick, domain=1.5:9.47, samples=2] {2 - 1.357143*(x - 8)};
  \node[eqpoint] at (axis cs:8,2) {};
  \node[anchor=north east, font=\footnotesize] at (axis cs:8,2) {$\omega$};
  \node[eqpoint] at (axis cs:4.737,6.429) {};
  \node[anchor=south west, font=\footnotesize] at (axis cs:4.8,6.5) {$x^*$};
  \node[anchor=west, font=\footnotesize, omDef] at (axis cs:6.6,4.0) {$u_A$};
  \node[anchor=east, font=\footnotesize, omThm] at (axis cs:2.6,8.6) {$u_B$};
  \node[anchor=east, font=\footnotesize] at (axis cs:8.9,0.5) {budget line};
\end{axis}
\node[anchor=north east, font=\footnotesize] at (current axis.north east) {$O_B$\,};
\node[anchor=south west, font=\footnotesize] at (current axis.south west) {\,$O_A$};
\end{tikzpicture}

{\small The \omterm{def:b2:exchange-economies:edgeworth-box}{Edgeworth box} of the running example. $A$'s origin is at the bottom left, $B$'s at the top
right. From the \omterm{def:b2:exchange-economies:endowment}{endowment} $\omega = (8, 2)$, the budget line at the equilibrium price $p = 19/14$ reaches
the \omterm{def:b2:exchange-economies:allocation}{allocation} $x^*$, where an indifference curve of each consumer touches it.}
\end{omfigure}

\section{Pareto efficiency and the contract curve}\label{sec:b2:exchange-economies:pareto}

\begin{definition}[Pareto-efficient allocation]\label{def:b2:exchange-economies:pareto-efficient}
A \omterm{def:b2:exchange-economies:feasible}{feasible allocation} $x$ is \emph{Pareto-efficient}\index{Pareto-efficient allocation} if no \omterm{def:b2:exchange-economies:feasible}{feasible
allocation} $x'$ makes every consumer at least as well off and at least one strictly better off:
$u_i(x'_i) \ge u_i(x_i)$ for all $i$ with a strict inequality for some $i$.
\end{definition}

Efficiency says nothing about fairness: the \omterm{def:b2:exchange-economies:allocation}{allocation} that gives $A$ everything is efficient, since any
change takes something from $A$. It is a test of waste, not a ranking of \omterm{def:b2:exchange-economies:allocation}{allocations}.

\begin{proposition}[Tangency]\label{prop:b2:exchange-economies:tangency}
Let $u_A$ and $u_B$ be differentiable, strictly increasing and strictly quasi-concave. An \omterm{def:b2:exchange-economies:allocation}{allocation} in
the interior of the \omterm{def:b2:exchange-economies:edgeworth-box}{Edgeworth box} is \omterm{def:b2:exchange-economies:pareto-efficient}{Pareto-efficient} if and only if $\mathit{MRS}_A = \mathit{MRS}_B$
there.
\end{proposition}

\begin{proof}
If $\mathit{MRS}_A > \mathit{MRS}_B$, $A$ values good 1 more, in units of good 2, than $B$ does. Let $B$
give $A$ a small amount $\varepsilon$ of good 1 against $t\varepsilon$ of good 2, with $\mathit{MRS}_B <
t < \mathit{MRS}_A$. To first order $A$'s utility changes by $\varepsilon\,\partial_2 u_A\,
(\mathit{MRS}_A - t) > 0$ and $B$'s by $\varepsilon\,\partial_2 u_B\,(t - \mathit{MRS}_B) > 0$: both gain,
so the \omterm{def:b2:exchange-economies:allocation}{allocation} is not efficient. The case $\mathit{MRS}_A < \mathit{MRS}_B$ is symmetric. Conversely,
if the rates are equal, the two indifference curves through the point are tangent; the set of bundles
$A$ weakly prefers is convex and lies on one side of the common tangent, and the set $B$ weakly prefers,
seen from the same corner, lies on the other side. Any point both weakly prefer lies on both sides, so it
is the point itself (strict quasi-concavity rules out a second common point): no \omterm{def:b2:exchange-economies:feasible}{feasible allocation}
improves on it.
\end{proof}

\begin{definition}[Contract curve]\label{def:b2:exchange-economies:contract-curve}
The \emph{contract curve}\index{contract curve} is the set of \omterm{def:b2:exchange-economies:pareto-efficient}{Pareto-efficient allocations} in the
\omterm{def:b2:exchange-economies:edgeworth-box}{Edgeworth box}.
\end{definition}

With Cobb--Douglas utilities the \omterm{def:b2:exchange-economies:contract-curve}{contract curve} has a closed form. $\mathit{MRS}_A = \frac{a}{1 - a}
\frac{x_{A2}}{x_{A1}}$ and $\mathit{MRS}_B = \frac{b}{1 - b}\frac{\bar\omega_2 - x_{A2}}{\bar\omega_1 -
x_{A1}}$; setting them equal and solving for $x_{A2}$,
\begin{equation}\label{eq:b2:exchange-economies:contract}
x_{A2} = \frac{b(1 - a)\,\bar\omega_2\,x_{A1}}{a(1 - b)(\bar\omega_1 - x_{A1}) + b(1 - a)\,x_{A1}} .
\end{equation}
It runs from one origin to the other. In the running example, $a = \tfrac12$ and $b = \tfrac23$, so $x_{A2}
= 20x_{A1}/(10 + x_{A1})$. When $a = b$, the denominator is $a(1 - a)\bar\omega_1$ and the curve is the
diagonal: consumers with the same shares consume goods in the same proportion at every efficient
\omterm{def:b2:exchange-economies:allocation}{allocation}.

At the \omterm{def:b2:exchange-economies:endowment}{endowment}, $\mathit{MRS}_A = 2/8 = 0.25$ and $\mathit{MRS}_B = 2 \times 8/2 = 8$: in units of good 2,
% code: exercises/solutions.py text_example
$B$ values good 1 at 32 times $A$'s valuation. There is a lot to trade. The \omterm{def:b2:exchange-economies:allocation}{allocations} both
consumers weakly prefer to the \omterm{def:b2:exchange-economies:endowment}{endowment} form the lens between their two indifference curves through
$\omega$, and the efficient ones among them are a piece of the \omterm{def:b2:exchange-economies:contract-curve}{contract curve}, from $x_{A1} = 3.26$ to
$x_{A1} = 6.27$.

\begin{omfigure}
\begin{tikzpicture}[econ]
\begin{axis}[omrecord, grid=none, axis lines=box, xmin=0, xmax=10, ymin=0, ymax=10, xtick={0,2,4,6,8,10},
    ytick={0,2,4,6,8,10}, xlabel={Good 1 held by $A$}, ylabel={Good 2 held by $A$},
    width=0.6\linewidth, height=0.6\linewidth, clip=true]
  \addplot[guide, domain=0:10, samples=2] {x};
  \addplot[black, very thick, domain=0:10, samples=60] {20*x/(10 + x)};
  \addplot[omDef, thick, domain=1.6:10, samples=60] {16/x};
  \addplot[omThm, thick, domain=0:8.2, samples=60] {10 - 32/(10 - x)^2};
  \addplot[omExo, line width=2.5pt, domain=3.257:6.266, samples=40] {20*x/(10 + x)};
  \node[eqpoint] at (axis cs:8,2) {};
  \node[anchor=north west, font=\footnotesize] at (axis cs:8,2) {$\omega$};
  \node[anchor=south, font=\footnotesize] (cc) at (axis cs:6.0,9.25) {contract curve};
  \draw[->, black!70] (cc.east) -- (axis cs:7.45,8.65);
  \node[anchor=south east, font=\footnotesize, omExo] at (axis cs:4.3,6.6) {core};
  \node[anchor=south west, font=\footnotesize, black!60] at (axis cs:1.2,1.6) {diagonal};
\end{axis}
\end{tikzpicture}

{\small The \omterm{def:b2:exchange-economies:contract-curve}{contract curve} of the running example, $x_{A2} = 20x_{A1}/(10 + x_{A1})$, lies above the
diagonal because $B$ spends more of his income on good 1 than $A$ does. The two indifference curves
through $\omega$ enclose the \omterm{def:b2:exchange-economies:allocation}{allocations} both prefer to the \omterm{def:b2:exchange-economies:endowment}{endowment}; the thick piece of the \omterm{def:b2:exchange-economies:contract-curve}{contract
curve} inside them is the \omterm{def:b2:exchange-economies:core}{core} of the economy (\cref{sec:b2:exchange-economies:several}).}
\end{omfigure}

\section{Competitive equilibrium}\label{sec:b2:exchange-economies:equilibrium}

Each consumer takes prices as given and sells her \omterm{def:b2:exchange-economies:endowment}{endowment} at those prices to buy her preferred bundle.
Her budget set is $\{x_i \ge 0 : p_1x_{i1} + p_2x_{i2} \le p_1\omega_{i1} + p_2\omega_{i2}\}$ and her
demand $x_i(p_1, p_2)$ is the best bundle in it.

\begin{definition}[Excess demand function]\label{def:b2:exchange-economies:excess-demand}
The \emph{excess demand function}\index{excess demand function} for good $\ell$ is $z_\ell(p_1, \dots,
p_n) = \sum_i x_{i\ell}(p_1, \dots, p_n) - \sum_i \omega_{i\ell}$: the total quantity demanded minus the
total quantity available.
\end{definition}

\begin{definition}[Competitive equilibrium]\label{def:b2:exchange-economies:competitive-equilibrium}
A \emph{competitive equilibrium}\index{competitive equilibrium (general)} of an \omterm{def:b2:exchange-economies:exchange-economy}{exchange economy} is a
price vector $p^* \gg 0$ and an \omterm{def:b2:exchange-economies:allocation}{allocation} $x^*$ such that each $x^*_i$ is consumer $i$'s demand at
$p^*$, and every market clears: $z_\ell(p^*) = 0$ for every good $\ell$.
\end{definition}

The Year 1 volume found the equilibrium of one market, taking the prices of all other goods as fixed.
Here every price is determined together, since a consumer's income is the value of her \omterm{def:b2:exchange-economies:endowment}{endowment} and
changes with every price.

\begin{definition}[Walras' law]\label{def:b2:exchange-economies:walras-law}
\emph{Walras' law}\index{Walras' law} is the identity $\sum_\ell p_\ell\,z_\ell(p) = 0$ for every price
vector $p$: the value of total excess demand is zero, whether or not markets clear.
\end{definition}

\begin{proposition}[Walras' law holds]\label{prop:b2:exchange-economies:walras}
If every consumer's preferences are locally non-satiated, so that each spends her whole income,
\omterm{def:b2:exchange-economies:walras-law}{Walras' law} holds. Hence if all markets but one clear at prices $p \gg 0$, the last clears too.
\end{proposition}

\begin{proof}
Each consumer's budget binds: $\sum_\ell p_\ell x_{i\ell}(p) = \sum_\ell p_\ell \omega_{i\ell}$, that is
$\sum_\ell p_\ell (x_{i\ell}(p) - \omega_{i\ell}) = 0$. Summing over $i$ gives $\sum_\ell p_\ell
z_\ell(p) = 0$. If $z_\ell(p) = 0$ for $\ell = 1, \dots, n - 1$, then $p_n z_n(p) = 0$ and, as $p_n > 0$,
$z_n(p) = 0$.
\end{proof}

Demands depend only on relative prices: multiplying every price by $t > 0$ multiplies each income by $t$
and leaves every budget set unchanged. So excess demand is homogeneous of degree zero, and one price can
be fixed.

\begin{definition}[Numeraire]\label{def:b2:exchange-economies:numeraire}
A \emph{numeraire}\index{numeraire} is a good whose price is set to one; every other price is then
measured in units of it.
\end{definition}

With $n$ goods there are $n - 1$ relative prices to find, and by \omterm{def:b2:exchange-economies:walras-law}{Walras' law} only $n - 1$ market-clearing
conditions to impose. With two goods and good 2 as \omterm{def:b2:exchange-economies:numeraire}{numeraire}, one equation in one unknown $p = p_1/p_2$:
$z_1(p) = 0$.

\begin{example}[The running example]\label{ex:b2:exchange-economies:running}
A Cobb--Douglas consumer spends the share $a$ of her income $p\omega_{A1} + \omega_{A2}$ on good 1, so
$x_{A1} = a(p\omega_{A1} + \omega_{A2})/p$. Market clearing for good 1, multiplied by $p$, reads
$a(p\omega_{A1} + \omega_{A2}) + b(p\omega_{B1} + \omega_{B2}) = p\bar\omega_1$, so
\[
p^* = \frac{a\,\omega_{A2} + b\,\omega_{B2}}{(1 - a)\,\omega_{A1} + (1 - b)\,\omega_{B1}}
= \frac{\tfrac12 \times 2 + \tfrac23 \times 8}{\tfrac12 \times 8 + \tfrac13 \times 2} =
\frac{19/3}{14/3} = \frac{19}{14} = 1.357 .
\]
% code: exercises/solutions.py text_example
$A$'s income is $1.357 \times 8 + 2 = 12.857$; she buys $x_{A1} = 0.5 \times 12.857/1.357 = 4.737$ and
$x_{A2} = 6.429$. $B$'s income is $1.357 \times 2 + 8 = 10.714$; he buys $x_{B1} = \tfrac23 \times
10.714/1.357 = 5.263$ and $x_{B2} = 3.571$. Both markets clear: $4.737 + 5.263 = 10$ and $6.429 + 3.571 =
10$, the second by \omterm{def:b2:exchange-economies:walras-law}{Walras' law}. The \omterm{def:b2:exchange-economies:allocation}{allocation} satisfies \eqref{eq:b2:exchange-economies:contract}:
$20 \times 4.737/14.737 = 6.429$.
\end{example}

The equilibrium lies on the \omterm{def:b2:exchange-economies:contract-curve}{contract curve}, as \cref{ch:b2:welfare-theorems} will prove in general: at
an interior equilibrium each consumer sets her $\mathit{MRS}$ equal to the same price ratio, so the two
rates are equal. It lies inside the lens, since each consumer could have kept her \omterm{def:b2:exchange-economies:endowment}{endowment}.

Equilibrium can also be read from the box with \emph{offer curves}. $A$'s offer curve is the set of her
demands as $p$ varies, drawn from her origin; $B$'s likewise. Every point of an offer curve is on a
budget line through $\omega$, and an equilibrium is a point where the two curves meet away from $\omega$:
there both consumers demand the same \omterm{def:b2:exchange-economies:allocation}{allocation} at the same price.

\begin{omfigure}
\begin{tikzpicture}[econ]
\begin{axis}[omrecord, grid=none, axis lines=box, xmin=0, xmax=10, ymin=0, ymax=10, xtick={0,2,4,6,8,10},
    ytick={0,2,4,6,8,10}, xlabel={Good 1 held by $A$}, ylabel={Good 2 held by $A$},
    width=0.6\linewidth, height=0.6\linewidth, clip=true]
  \addplot[omDef, thick, domain=0.2:8, samples=120, variable=\t] ({0.5*(8*\t + 2)/\t}, {0.5*(8*\t + 2)});
  \addplot[omThm, thick, domain=0.12:8, samples=120, variable=\t]
    ({10 - (2/3)*(2*\t + 8)/\t}, {10 - (1/3)*(2*\t + 8)});
  \addplot[black, domain=2.7:9.47, samples=2] {2 - 1.357143*(x - 8)};
  \node[eqpoint] at (axis cs:8,2) {};
  \node[anchor=north east, font=\footnotesize] at (axis cs:8,2) {$\omega$};
  \node[eqpoint] at (axis cs:4.737,6.429) {};
  \node[anchor=south west, font=\footnotesize] at (axis cs:4.8,6.5) {$x^*$};
  \node[anchor=west, font=\footnotesize, omDef] at (axis cs:5.0,9.0) {$A$'s offer curve};
  \node[anchor=south west, font=\footnotesize, omThm] at (axis cs:0.4,7.2) {$B$'s offer curve};
\end{axis}
\end{tikzpicture}

{\small Offer curves of the running example. As $p$ rises, $A$, who owns good 1, becomes richer and
moves up her curve; $B$'s curve is traced from his corner. Both curves pass through $\omega$, but at
different prices ($A$ chooses not to trade at $p = 0.25$, $B$ at $p = 8$), so that meeting is not an
equilibrium. They meet again at $x^*$, on the budget line of slope $-19/14$, where both demand the same
\omterm{def:b2:exchange-economies:allocation}{allocation} at the same price.}
\end{omfigure}

\begin{record}[Markets in the laboratory]
Smith (1962) ran ten classroom markets of this kind, in which buyers and sellers held private
% ledger: b2:smith1962-t6a, b2:smith1962-t7
reservation prices and traded by open outcry. He measured each period's dispersion around the predicted
price by a \emph{coefficient of convergence}, $\alpha = 100\,\sigma_0/P_0$, where $\sigma_0$ is the
root mean square deviation of contract prices from the predicted price $P_0$. In test 6A, $\alpha$ fell
from 53.8 in the first period to 9.4 in the fourth; in test 7, from 49.1 to 2.7 over six periods. In all
% figdata: smith_summary (share_alpha_fell 100)
ten tests with three or more periods it was lower in the last period than in the first. Quantities
% ledger: b2:smith1962-t1
traded were close to the prediction from the start: five, five, five, seven and six units against a
predicted six in test 1. Each market was run within one class period, with the students of one class
as its traders: far from the large numbers that justify price taking in theory.
\end{record}

\section{Existence}\label{sec:b2:exchange-economies:existence}

A \omterm{def:b2:exchange-economies:competitive-equilibrium}{competitive equilibrium} need not exist: if a consumer's demand jumps as prices change, excess demand
can jump over zero. Two conditions are enough.

\begin{theorem}[Existence of a competitive equilibrium]\label{thm:b2:exchange-economies:existence}
Suppose every consumer has a positive \omterm{def:b2:exchange-economies:endowment}{endowment} of every good, and continuous, strictly increasing and
strictly quasi-concave preferences. Then the \omterm{def:b2:exchange-economies:exchange-economy}{exchange economy} has a \omterm{def:b2:exchange-economies:competitive-equilibrium}{competitive equilibrium}.
\end{theorem}

\admitted

The general proof uses a fixed-point theorem: the map that raises the prices of goods in excess demand
and lowers the others, normalised so that prices sum to one, is continuous on the simplex and has a
fixed point, at which no market is in excess demand. With two goods the intermediate value theorem is
enough, and the two conditions of the theorem show where they are used.

\begin{proposition}[Two goods]\label{prop:b2:exchange-economies:two-goods}
With two goods and the assumptions of \cref{thm:b2:exchange-economies:existence}, $z_1(p)$ is continuous on
$(0, \infty)$, positive for $p$ small enough and negative for $p$ large enough. Hence it is zero at some
$p^* > 0$, and by \omterm{def:b2:exchange-economies:walras-law}{Walras' law} both markets clear there.
\end{proposition}

\begin{proof}
Continuity of preferences and strict quasi-concavity make each demand a continuous function of $p$
(admitted from the theory of the consumer). As $p \to 0$, good 1 becomes free while every consumer's
income stays at least $\omega_{i2} > 0$; a consumer with strictly increasing preferences then demands
more good 1 than any bound, and $z_1(p) \to +\infty$. By the same argument with the roles of the goods
swapped, $z_2(p) \to +\infty$ as $p \to \infty$, and \omterm{def:b2:exchange-economies:walras-law}{Walras' law} $pz_1 + z_2 = 0$ gives $z_1(p) < 0$ for
$p$ large. The intermediate value theorem gives a zero.
\end{proof}

The conditions matter. Without continuity, demand can jump past the total \omterm{def:b2:exchange-economies:endowment}{endowment}. Without positive
\omterm{def:b2:exchange-economies:endowment}{endowments}, a consumer who owns only good 1 has no income when $p \to 0$, and her demand need not grow:
excess demand can stay negative all the way down.

\section{Several equilibria and stability}\label{sec:b2:exchange-economies:several}

Nothing in the assumptions makes the equilibrium unique. Take $u_A = x_{A1} - \tfrac18 x_{A2}^{-8}$ and
$u_B = -\tfrac18 x_{B1}^{-8} + x_{B2}$, with $\omega_A = (2, r)$ and $\omega_B = (r, 2)$, where $r =
% code: exercises/solutions.py three
2^{8/9} - 2^{1/9} = 0.7717$. $A$ cares about good 2 only up to a point, $B$ about good 1 only up to a
point. The tangency conditions give $x_{A2} = p^{1/9}$ and $x_{B1} = p^{-1/9}$; $A$ spends the rest of
her income on good 1, $x_{A1} = 2 + (r - p^{1/9})/p$. So
\[
z_1(p) = \frac{r - p^{1/9}}{p} + p^{-1/9} - r ,
\]
and $z_1(\tfrac12) = z_1(1) = z_1(2) = 0$ (\cref{exo:b2:exchange-economies:11}): three equilibria.

\begin{omfigure}
\begin{tikzpicture}[econ]
\begin{axis}[omrecord, grid=none, xmode=log, log ticks with fixed point, xmin=0.35, xmax=3.5,
    ymin=-7, ymax=12, xtick={0.5,1,2}, ytick={-5,0,5,10}, xlabel={Relative price $p = p_1/p_2$
    (log scale)}, ylabel={$1000 \times z_1(p)$}, width=0.62\linewidth, height=0.45\linewidth]
  \addplot[guide, domain=0.35:3.5, samples=2] {0};
  \addplot[omDef, very thick, domain=0.35:3.5, samples=150]
    {1000*((0.7716897 - x^(1/9))/x + x^(-1/9) - 0.7716897)};
  \node[eqpoint] at (axis cs:0.5,0) {};
  \node[eqpoint] at (axis cs:1,0) {};
  \node[eqpoint] at (axis cs:2,0) {};
  \node[anchor=north east, font=\footnotesize] at (axis cs:0.49,-0.6) {stable};
  \node[anchor=south, font=\footnotesize] at (axis cs:1,0.6) {unstable};
  \node[anchor=south, font=\footnotesize] at (axis cs:2,0.6) {stable};
\end{axis}
\end{tikzpicture}

{\small Excess demand for good 1 in an economy with three equilibria, $p = \tfrac12$, $1$ and $2$ (in
thousandths of a unit, since the swings are small). Where $z_1$ crosses zero from above, a price that
rises with excess demand returns to the equilibrium; where it crosses from below, it moves away.}
\end{omfigure}

Which equilibrium would a market reach? Imagine an auctioneer who calls out prices, collects demands
without letting anyone trade, and raises the price of a good in excess demand: \emph{tâtonnement}
(groping), $\dot p = k\,z_1(p)$ with $k > 0$. An equilibrium is \emph{locally stable} under this process if
prices that start near it converge to it, which with two goods happens exactly when $z_1$ is decreasing
there. In the example, $p = \tfrac12$ and $p = 2$ are stable and $p = 1$ is not. The listing runs the
process from seven starting prices.

\omcode{code/bachelor-2/06-exchange-economies/listings/tatonnement.py}{1}{21}{Tâtonnement from several
starting prices in the economy with three equilibria.}\label{lst:b2:exchange-economies:tatonnement}

\omoutput{code/bachelor-2/06-exchange-economies/outputs/tatonnement.txt}{Prices that start below $1$ go
to $\tfrac12$, prices that start above go to $2$.}\label{out:b2:exchange-economies:tatonnement}

% code: listings/tatonnement.py
A start at exactly 1 stays there; a start at 0.99 goes to $\tfrac12$ and one at 1.01 to $2$. Two
cautions. Tâtonnement is a thought experiment: in real markets trade happens at prices out of
equilibrium, which changes \omterm{def:b2:exchange-economies:endowment}{endowments} and so the equilibrium itself. And with three goods or more,
examples are known in which tâtonnement cycles for ever around a unique equilibrium, so stability is not
guaranteed by the assumptions that guarantee existence. The laboratory markets of the opening chart did
converge, but they were single markets with supply and demand designed by the experimenter.

With many equilibria, the theory predicts a set, and which point the economy reaches depends on history.
With few traders it predicts less still: two consumers facing each other in a box have no reason to take
prices as given, and might bargain to any efficient point both prefer to their \omterm{def:b2:exchange-economies:endowment}{endowment}.

\begin{definition}[Core]\label{def:b2:exchange-economies:core}
A coalition $S$ of consumers \emph{blocks} a \omterm{def:b2:exchange-economies:feasible}{feasible allocation} $x$ if its members, trading only their
own \omterm{def:b2:exchange-economies:endowment}{endowments} among themselves, can reach bundles $(y_i)_{i \in S}$ with $\sum_{i \in S} y_i \le
\sum_{i \in S}\omega_i$ that make none of them worse off than under $x$ and at least one strictly better
off. The \emph{core}\index{core} is the set of \omterm{def:b2:exchange-economies:feasible}{feasible allocations} that no coalition blocks.
\end{definition}

With two consumers, a single consumer blocks any \omterm{def:b2:exchange-economies:allocation}{allocation} she likes less than her \omterm{def:b2:exchange-economies:endowment}{endowment}, and the
pair blocks any inefficient \omterm{def:b2:exchange-economies:allocation}{allocation}. So the \omterm{def:b2:exchange-economies:core}{core} is the part of the \omterm{def:b2:exchange-economies:contract-curve}{contract curve} inside the lens,
the thick piece in the second figure. Every \omterm{def:b2:exchange-economies:competitive-equilibrium}{competitive equilibrium} is in the \omterm{def:b2:exchange-economies:core}{core}: a coalition that
could do better with its own \omterm{def:b2:exchange-economies:endowment}{endowments} could have bought the better bundles at equilibrium prices,
which contradicts the consumers' choices (\cref{ch:b2:welfare-theorems}). As the economy is replicated,
with many copies of each type of consumer, more coalitions can form and the \omterm{def:b2:exchange-economies:core}{core} shrinks; in the limit
only competitive \omterm{def:b2:exchange-economies:allocation}{allocations} remain. \Cref{exo:b2:exchange-economies:12} proves the first step. Price
taking is then a conclusion about large economies, not an assumption about every economy.

\begin{model}{Walrasian exchange}
\assumes Consumers take prices as given; preferences are convex and monotone; there are no externalities;
there is a market for every good.
\predicts Prices that clear all markets together; equilibrium \omterm{def:b2:exchange-economies:allocation}{allocations} on the \omterm{def:b2:exchange-economies:contract-curve}{contract curve} and in
the \omterm{def:b2:exchange-economies:core}{core}; changes in relative prices when \omterm{def:b2:exchange-economies:endowment}{endowments} are redistributed between consumers with different
tastes.
\tested Laboratory double-auction markets, in which prices converge towards the competitive prediction
within a few periods (Smith 1962 and the experiments that followed).
\breaks With few traders, price taking has no justification and the outcome is a bargain anywhere in the
\omterm{def:b2:exchange-economies:core}{core}; with several equilibria, the theory does not say which one is reached; with three goods or more,
price adjustment need not converge; and the model is silent on how trade proceeds out of equilibrium.
\end{model}

\section{Exercises}\label{sec:b2:exchange-economies:exercises}

\begin{exercise}[$\star$]\label{exo:b2:exchange-economies:1}
The total \omterm{def:b2:exchange-economies:endowment}{endowment} is $(10, 10)$. Which of these \omterm{def:b2:exchange-economies:allocation}{allocations} $(x_A, x_B)$ are \omterm{def:b2:exchange-economies:feasible}{feasible}, and which
non-wasteful? (i) $(4, 7), (6, 3)$; (ii) $(5, 5), (6, 4)$; (iii) $(3, 2), (6, 7)$.
\end{exercise}

\begin{exercise}[$\star$]\label{exo:b2:exchange-economies:2}
$u_A = x_{A1}x_{A2}$ and $u_B = x_{B1}x_{B2}^2$. At the \omterm{def:b2:exchange-economies:allocation}{allocation} $x_A = (4, 6)$, $x_B = (6, 4)$, compute
both marginal rates of substitution. Is the \omterm{def:b2:exchange-economies:allocation}{allocation} \omterm{def:b2:exchange-economies:pareto-efficient}{Pareto-efficient}? If not, which way should good 1
move?
\end{exercise}

\begin{exercise}[$\star$]\label{exo:b2:exchange-economies:3}
Read the opening chart. In which test was the first period's price furthest from the prediction, and
in which tests was the last period's price still more than 5 per cent away? Does any line cross the
prediction?
\end{exercise}

\begin{exercise}[$\star$]\label{exo:b2:exchange-economies:4}
Describe a laboratory result that would refute the claim that double-auction markets converge to the
competitive prediction, and say why a single period far from the prediction would not.
\end{exercise}

\begin{exercise}[$\star\star$]\label{exo:b2:exchange-economies:5}
$u_A = \ln x_{A1} + \ln x_{A2}$ with $\omega_A = (6, 2)$, and $u_B = \ln x_{B1} + 3\ln x_{B2}$ with
$\omega_B = (4, 8)$. Find the \omterm{def:b2:exchange-economies:competitive-equilibrium}{competitive equilibrium} price $p = p_1/p_2$ and \omterm{def:b2:exchange-economies:allocation}{allocation}, and check that
both markets clear.
\end{exercise}

\begin{exercise}[$\star\star$]\label{exo:b2:exchange-economies:6}
For the economy of \cref{exo:b2:exchange-economies:5}, write the \omterm{def:b2:exchange-economies:contract-curve}{contract curve}, compute $x_{A2}$ at
$x_{A1} = 2$, $5$ and $8$, and check that the equilibrium \omterm{def:b2:exchange-economies:allocation}{allocation} lies on it. Is it above or below the
diagonal, and why?
\end{exercise}

\begin{exercise}[$\star\star$]\label{exo:b2:exchange-economies:7}
In the economy of \cref{exo:b2:exchange-economies:5}, $B$ gives $A$ two units of good 2 before trade
opens. Predict the direction of the change in $p$ before computing it, then compute the new price. Show
that with Cobb--Douglas consumers a transfer of $\Delta$ units of good 2 from $B$ to $A$ changes $p^*$ by
$(a - b)\Delta/((1 - a)\omega_{A1} + (1 - b)\omega_{B1})$.
\end{exercise}

\begin{exercise}[$\star\star$]\label{exo:b2:exchange-economies:8}
In the economy of \cref{exo:b2:exchange-economies:5}, compute $z_1$ and $z_2$ at $p = 1$, and check
\omterm{def:b2:exchange-economies:walras-law}{Walras' law}. Which price must rise?
\end{exercise}

\begin{exercise}[$\star\star$]\label{exo:b2:exchange-economies:9}
Design a laboratory test of whether a market with two competitive equilibria settles at one of them
predictably. What would you vary across sessions, and what result would show that history, not the
model, picks the equilibrium?
\end{exercise}

\begin{exercise}[$\star\star\star$]\label{exo:b2:exchange-economies:10}
Prove \omterm{def:b2:exchange-economies:walras-law}{Walras' law} for $I$ consumers and $n$ goods, with locally non-satiated preferences. Show that if
$z_\ell(p) = 0$ for $n - 1$ goods and $p_n > 0$, the $n$-th market clears, and give an example in which
a good with a zero price is in excess supply at equilibrium.
\end{exercise}

\begin{exercise}[$\star\star\star$]\label{exo:b2:exchange-economies:11}
For the economy of \cref{sec:b2:exchange-economies:several}: (a) derive $x_{A2} = p^{1/9}$, $x_{B1} =
p^{-1/9}$ and $z_1(p)$; (b) check that $p = \tfrac12$, $1$ and $2$ are equilibria; (c) compute the sign
of $z_1'$ at each and say which are stable under tâtonnement.
\end{exercise}

\begin{exercise}[$\star\star\star$]\label{exo:b2:exchange-economies:12}
Replicate a two-consumer economy: two consumers of type $A$ with \omterm{def:b2:exchange-economies:endowment}{endowment} $\omega_A$ and two of type
$B$ with $\omega_B$, all with strictly quasi-concave preferences of their type. Let $(x_A, x_B)$ be an
\omterm{def:b2:exchange-economies:allocation}{allocation} of the two-consumer economy on its \omterm{def:b2:exchange-economies:contract-curve}{contract curve} with $x_A \neq \omega_A$ and $u_A(x_A) =
u_A(\omega_A)$, $u_B(x_B) > u_B(\omega_B)$, and give each consumer of the replica her type's bundle.
Show that the coalition of the two $A$ consumers and one $B$ consumer blocks this \omterm{def:b2:exchange-economies:allocation}{allocation}, so that it
is in the \omterm{def:b2:exchange-economies:core}{core} of the two-consumer economy but not of its replica.
\end{exercise}

\section{Problem: Fish for grain}\label{sec:b2:exchange-economies:problem}

\begin{problem}[{Weekend problem --- Brenland's fishers and Kestria's farmers, in two invented
economies, trade fish for grain, find the price that clears both markets, and see what a gift of grain
does to it}]
\label{pb:b2:exchange-economies:1}
Treat each country as one consumer. Brenland holds 60 tonnes of fish and 10 of grain, Kestria 20 of fish
and 90 of grain. Brenland's utility is $\tfrac12\ln f + \tfrac12\ln g$, Kestria's $\tfrac14\ln f +
\tfrac34\ln g$, where $f$ and $g$ are tonnes of fish and grain. Grain is the \omterm{def:b2:exchange-economies:numeraire}{numeraire} and $p$ is the
price of a tonne of fish in tonnes of grain. Round prices and marginal rates of substitution to four
decimals and quantities to two.

\medskip
\noindent\textbf{Part I --- The box.}
\begin{enumerate}
  \item Give the dimensions of the \omterm{def:b2:exchange-economies:edgeworth-box}{Edgeworth box} and the \omterm{def:b2:exchange-economies:endowment}{endowment} point, with Brenland at the
        bottom-left corner.
  \item Compute each country's marginal rate of substitution of grain for fish at the \omterm{def:b2:exchange-economies:endowment}{endowment}.
  \item Is the \omterm{def:b2:exchange-economies:endowment}{endowment} \omterm{def:b2:exchange-economies:pareto-efficient}{Pareto-efficient}? Within which range of terms of trade (grain per tonne of fish)
        would both gain from a small trade?
  \item Brenland offers 10 tonnes of fish for 5 tonnes of grain. Check that both gain.
  \item Who exports fish, and why, in one sentence?
\end{enumerate}

\medskip
\noindent\textbf{Part II --- The Pareto set.}
\begin{enumerate}[resume]
  \item Write the condition for an interior \omterm{def:b2:exchange-economies:allocation}{allocation} to be \omterm{def:b2:exchange-economies:pareto-efficient}{Pareto-efficient}.
  \item Derive the \omterm{def:b2:exchange-economies:contract-curve}{contract curve} as Brenland's grain $g_B$ as a function of its fish $f_B$.
  \item Check that it passes through both corners of the box. Is it above or below the diagonal?
  \item Find the range of $f_B$ over which the \omterm{def:b2:exchange-economies:contract-curve}{contract curve} lies inside the lens of \omterm{def:b2:exchange-economies:allocation}{allocations} both
        countries weakly prefer to the \omterm{def:b2:exchange-economies:endowment}{endowment}.
  \item Is the \omterm{def:b2:exchange-economies:allocation}{allocation} that gives Brenland all the fish and all the grain \omterm{def:b2:exchange-economies:pareto-efficient}{Pareto-efficient}? What does
        your answer say about efficiency?
\end{enumerate}

\medskip
\noindent\textbf{Part III --- Equilibrium.}
\begin{enumerate}[resume]
  \item Write each country's demand for fish and for grain as a function of $p$.
  \item Write the excess demand for fish $z_f(p)$.
  \item Find the equilibrium price.
  \item Find the equilibrium \omterm{def:b2:exchange-economies:allocation}{allocation} and each country's trade.
  \item Check that the grain market clears, and name the law that says it must.
  \item Check that the \omterm{def:b2:exchange-economies:allocation}{allocation} lies on the \omterm{def:b2:exchange-economies:contract-curve}{contract curve} and inside the range of question 9.
  \item Compare each country's utility (or $f^{a}g^{1-a}$) with its \omterm{def:b2:exchange-economies:endowment}{endowment}'s.
  \item For each country, find the extra grain that, added to its \omterm{def:b2:exchange-economies:endowment}{endowment} without trade, would make it
        as well off as trade does.
\end{enumerate}

\medskip
\noindent\textbf{Part IV --- A gift of grain.} Before trade opens, Kestria gives Brenland 10 tonnes of
grain.
\begin{enumerate}[resume]
  \item Predict the direction of the change in the price of fish, and explain it.
  \item Compute the new equilibrium price and \omterm{def:b2:exchange-economies:allocation}{allocation}.
  \item Who gains and who loses relative to Part III? Does Kestria still gain from trade, given its new
        \omterm{def:b2:exchange-economies:endowment}{endowment}?
  \item Check that the new \omterm{def:b2:exchange-economies:allocation}{allocation} lies on the same \omterm{def:b2:exchange-economies:contract-curve}{contract curve}.
  \item Could a transfer of \omterm{def:b2:exchange-economies:endowment}{endowments} reach any point of the \omterm{def:b2:exchange-economies:contract-curve}{contract curve}? What does that suggest
        for the next chapter?
  \item State the result: the equilibrium price of fish in grain, and each country's gain from trade as
        an equivalent amount of grain.
\end{enumerate}
\end{problem}
